\section{Why five missing counts can be recovered}
The neighbour counts at any 507 of the 512 vertices of $Q_9$ determine
the selected vertices uniquely. If the selection lies on a known
parity side, any 251 of the 256 counts on the opposite side suffice.
This extends the preceding four-count result.

Subtract two selections that agree on the known counts. The difference
$f$ has entries $-1,0,1$, and its neighbour sums can be nonzero only
at the five missing positions. If the positions have mixed parity,
split $f$ by parity: each part has at most four missing counts, already
covered. If the five positions contain no pair at distance eight,
the preceding eight-count result applies.

The remaining case has five positions of one parity and a pair at
distance eight. Move one endpoint of that pair to zero and the other
to $(1,1,1,1,1,1,1,1,0)$. The remaining three positions give one of eight
three-bit patterns in each of the first eight coordinates. Count the
occurrences of each pattern; even parity fixes their ninth bits.
This describes every possible configuration with 6435 finite choices.
Discarding repeated positions and identifying equal distance matrices
leaves 372 matrices to consider~\cite{research240}.

Apply the inverse formula at the five antipodes. The unknown missing
count differences $w$ satisfy
\[
Mw=315s,\qquad w\in\mathbb Z^5,\qquad s\in\{-1,0,1\}^5.
\]
The diagonal of $M$ is 128, and the other coefficients are
$-16,8,-10,35$ at distances $2,4,6,8$. Multiplication by its adjugate
$B$ gives $Dw_i=315(Bs)_i$, where $D>0$ is its determinant. Hence
$q=D/\gcd(D,315)$ must divide every $(Bs)_i$. If
\[
q>\sum_j|B_{ij}|,
\]
then the permitted values of $s$ make $(Bs)_i$ too small to be a nonzero
multiple of $q$, forcing $w_i=0$.

This inequality settles all five unknowns for 334 matrices. It settles
four unknowns for another five matrices; the last equation becomes
$128w_i=315s_i$, which vanishes by parity. For the remaining 33 matrices,
the 243 possible ternary vectors can be checked directly: only zero
satisfies all five divisibility conditions. Thus every missing count
difference vanishes, and the inverse gives $f=0$ everywhere.

The saved native calculation covers all 90396 matrix/right-hand-side
cases using exact integers. The 372 corresponding Lean statements
match the saved matrices and all passed Lean checking. A separate
enumeration of all 2699004 triples of actual vertices outside the fixed
pair recovers the same 372 matrix types. The passage from arbitrary cube positions to those
matrices and the parity reduction remain written arguments with
native computational support. The complete actual-cube Lean theorem
still covers three missing counts.

Two different selections must therefore change at least six neighbour
counts. A single-vertex change alters nine, leaving the minimum between
six and nine. A failure with six, seven, or eight missing counts would
have to contain a pair at distance eight. This reconstruction result
does not supply a longer snake or change the global length bound.
